\section{Human--AI Research Protocol}\label{app:protocol}

This appendix documents how the paper was produced. arXiv's policy is that generative AI tools cannot be authors, that authors must report significant use of such tools, and that authors take full responsibility for all content however it was generated \citep{arxiv2023policy,arxiv2026policy}. We give more than a disclosure statement because the structure of the collaboration bears on how to judge the work. The first two versions were produced on 28 and 29 October 2026; the second review round ran from 29 October to 2 November (six of its eight reviews stopped on 29 October at the account's usage limit and resumed from saved state on 2 November), and this third version was revised on 2 November and re-estimated and revised again on 3 November. Times are those of the session clock (UTC$-$7 in October, UTC$-$8 from 1 November), checked against the platform's UTC timestamps. Section~\ref{supp:protocol} of the Supplementary Appendix gives the detailed records.

\subsection{Agents and models}\label{app:protocol-agents}
The author directed the research and made its decisions; AI agents drafted the text, code and proofs. The author posed the research questions, set the scope and made every decision listed in Table~\ref{tab:decisions}, and is responsible for the paper. She read the manuscript and commissioned four independent reviews of its third version by models of other developers, run outside this workflow (Section~\ref{supp:protocol-outside} of the Supplementary Appendix); their findings were verified item by item and answered in the response to the outside reviews. All AI work was done by agents built on Anthropic's Claude models inside one research session. The session was configured for the model identifier \texttt{claude-opus-5-5}, which agents inherited unless another model was chosen for them. After the author asked for different models to be used (13:08 on 29 October), the cross-model referees of the first round and the second assumptions auditor ran on \texttt{claude-fable-5-1} and the literature agent of the second version on \texttt{claude-sonnet-5-5}; in the second review round the four referees ran on \texttt{claude-fable-5-1} and the four practitioner panels on \texttt{claude-sonnet-5-5}. A \emph{lead agent} coordinated the work and was the only agent that communicated with the author. Where independence mattered, an agent was given a distinct role, was not given the others' outputs and was told not to read them; all agents worked in one shared file system. Every agent had the same tools: reading and writing files in the shared project folder, running code, and web search; bulk SEC data were retrieved through a browser on the author's computer, with her approval, on 28 and 29 October, twice on 2 November and twice on 3 November (Section~\ref{app:data} of the Supplementary Appendix). Roles were set by instruction, not by permissions, so reviewers were told not to edit the manuscript, the code or each other's reports. The session transcripts are not part of the replication package; the agent log, every review report, the issue logs and the responses to reviewers are. The agents fell into five groups, drafting and integration, theory, data and parameters, literature and review, which Section~\ref{supp:protocol} of the Supplementary Appendix describes (Table~\ref{tab:agents}). The project later grew into three papers: this one, a companion paper on the economics of AI transformation \citep{murthy2026transformation}, and a paper on the software engine.

\subsection{Review rounds}\label{app:protocol-rounds}
A council of three agents critiqued the lead agent's first answer to the research question without seeing each other's critiques; four of their changes survive here (capital intensity in the frontier, very high scores at large scale during market transitions, the measurement basis and accountability costs that scale with the damage a failure can cause). A referee panel of a theorist, a financial economist and an econometrician, and an audit agent, then reviewed the complete first draft; their reports led to the rebuilt panel as first reported, sample-size-invariant statistics, the treatment of $\alpha$ as a calibration range, the dynamic weight of \autoref{prop:dynamic} and the firm-by-firm inversion of Section~\ref{sec:horizons}. The first complete version was then reviewed on a frozen snapshot, taken at 12:46 on 29 October, by twelve panels on the models named in Section~\ref{app:protocol-agents}: four academic referees, four practitioner panels and four cross-model referees in the same fields, told not to read the other reports, with a second assumptions auditor. Each practitioner panel was a single agent playing several named roles, so the panels are structured self-critique, not a set of independent human experts. The eight referee reports recommended major revision, the four practitioner panels asked for changes before release, and no referee found a false theorem in the static model (Section~\ref{app:protocol-panels} of the Supplementary Appendix gives their findings). The lead agent consolidated the findings into a revision plan, on which the author took four decisions (D-1 to D-4, Table~\ref{tab:decisions}). In the second version a sample-audit agent rebuilt the sample under five rules and re-estimated every statistic, a parameters agent adopted the re-grounded parameters, a literature agent corrected the citations, theory agents addressed the theory referees' findings, and the editor-author agent revised the text and split the paper into this main text and a Supplementary Appendix.

The second version was reviewed on a frozen snapshot taken at 18:42 on 29 October by eight panels, told not to read one another's reports: four referees on \texttt{claude-fable-5-1} (economic theory; econometrics; finance and accounting; industrial organization and the economics of AI) and four practitioner panels on \texttt{claude-sonnet-5-5} (corporate executives; investors and bankers; investment management; risk, audit and regulation). The referees recommended minor revision; the panels asked for more, mainly on the empirical headline and the material for practitioners. The lead agent's revision plan of 2 November listed the findings; the author's decisions on it and on the work that followed are the six rows of Table~\ref{tab:decisions} dated 2 November. In the third version a coding agent recomputed the numbers the reviewers questioned, the lead agent retrieved further SEC data, and the editor-author agent revised the text; the round-2 response records the disposition of every finding. The author then commissioned three reviews of the third version from models of other developers, run outside this workflow: DeepSeek V4 Pro (high effort), GPT-6.1 Sol (extra-high effort) and GLM 5.3 (high effort). Their findings were verified item by item in this session, the author adopted every fix the verification supported (Table~\ref{tab:decisions}), and the response to the outside reviews records the disposition of each finding. A joint-parameter stress test of the calibrated claims was added in the same round (Section~\ref{supp:stress} of the Supplementary Appendix). An agent on \texttt{claude-fable-5-1} that had not taken part in the revision then audited it against the three issue logs, the author's decisions and the code outputs. Its one major finding was that the stress test moved the compute share over a narrower range than the scenario range of Section~\ref{sec:horizons-ceiling}; checking it showed that the floor of that range had been set in the wrong units (a cost per dollar of labor used as a share of revenue), which was corrected, and the stress test was rerun over the corrected range with a second seed; at the author's request it was then repeated with \StDraws{} draws per seed, run on her computer (Section~\ref{supp:stress} of the Supplementary Appendix). On 3 November the author added a fourth outside review, Mistral Large 4 (high effort), verified the same way, and five tests: an audit of a seeded sample of firm-years against the filings, which led to the correction of capitalized software read from the filings; business descriptions of every core-sample registrant, so that rule R2 applies to the main sample as defined, with two independent readings as a sensitivity; a specification curve of the empirical claims (Section~\ref{supp:speccurve} of the Supplementary Appendix); a reproduction of the whole pipeline from a fresh clone in a fresh environment, which matched the committed data, tables and main paper and found two stale records and figures whose layout depended on the Python environment, all regenerated; and a check of the legal and tax citations at their primary sources. On 4 November the author added three analyses that had been built for a plain-language companion page and checked there by an agent that had not built them (Section~\ref{supp:extensions} of the Supplementary Appendix). All four outside reviewers then re-reviewed the build of 4 November: GPT-6.1 Sol (extra-high effort) asked for a targeted major revision, GLM 5.3 (high effort) for release as a working paper after minor fixes, and Mistral Large 4 and DeepSeek V4 Pro (both high effort) for release now. Two agents on \texttt{claude-fable-5-1} verified the findings of the first two item by item and the lead agent verified the other two (logs in \texttt{notes/reviews/paper1\_v4/author\_reviews/}); no finding moved a headline number. On 5 November the author adopted the fixes the verification of the first two supported (Table~\ref{tab:decisions}), an agent on \texttt{claude-fable-5-1} that had not made them audited them, and the DeepSeek re-review added one clause to the introduction, which the author kept: that (H) is checked numerically and not proved.

\subsection{Decisions made by the author}
{\footnotesize
\begin{longtable}{@{}p{0.30\textwidth}p{0.66\textwidth}@{}}
\caption{Decisions made by the human author}\label{tab:decisions}\\
\toprule
Decision & Choice\\\midrule
\endfirsthead
\multicolumn{2}{@{}l}{\emph{Table \thetable{} continued}}\\\toprule
Decision & Choice\\\midrule
\endhead
\midrule\multicolumn{2}{r@{}}{\emph{Continued on the next page}}\\
\endfoot
\bottomrule
\endlastfoot
Research question and horizons & Limit of the Rule of 40 and successors under full automation, bounded to two, five and ten years, with necessary conditions\\
Independent reviewers & Requested a council of agents before the paper, a research team for the paper, and referee and practitioner panels\\
Evidence base & Theory plus firm-level SEC data, rather than theory with simulations only or a proprietary dataset\\
Authorship and affiliation & Sole human author, independent; AI collaboration documented as a method in the paper itself\\
Data access & Granted folder and browser access on the author's computer for SEC retrieval\\
Standard of evidence & ``Accurate, factual, referenced, cited and complete''; every assumption grounded in solid precedent; recent (AI-era) theory\\
Scope of the frontier & Apply every constraint of the transformation model to the frontier itself; account for agentic misbehavior, errors and hallucination\\
Scope of internal control & What matters is the processes, systems, integrations and policies that affect the timeliness, accuracy and completeness of financial reporting\\
Publication (10:50) & Three papers; release this one as soon as it is ready\\
Competing interests (10:50; 2 November) & No statement at first; on 2 November a minimal statement, drafted by the lead agent and confirmed by the author at 10:45: no funding and no competing interests\\
Estimands (11:19) & Asked what rigor requires on reporting results with and without acquisition-affected years\\
Practical outcome (11:40, 11:58) & Strategy by Rule-of-$X$ level and organization setup, down to the process step; this paper carries a short section, the companion paper the full playbook\\
Title (12:03) & ``The Rule of 40 as a Price Line: Efficient Frontiers, Economic Profit and the Limits of AI Automation''\\
Review independence (13:08) & Asked for different models to be used; the cross-model referees ran on a second model\\
Revision decisions (13:53) & Withdraw the 7-to-13 growth weight and report bounds (D-1); on the frontier object, asked for the lead agent's call (D-2); claims must hold under every view of the sample (D-3); adopt the re-grounded parameters in one commit (D-4)\\
Data for the sample rules (16:05) & Approved the download of the SEC data that rules R2 to R5 and the score before interest and taxes need, through the desktop application's built-in browser\\
Authorship statement (2 November) & ``The author directed the research and made its decisions; AI agents drafted the text, code and proofs''\\
Author's review (2 November) & Read the manuscript and commissioned three outside reviews of the third version (DeepSeek V4 Pro, GPT-6.1 Sol, GLM 5.3), run outside this workflow\\
SEC data for the third version (2 November) & Approved the retrieval of interest income, deferred revenue and contract liabilities for the core firms and of the acquirer's reports and prospectus cash flows, through the built-in browser, with nothing entered on any site\\
Basis of selling and marketing (2 November, 10:15, 12:35) & Net of maintenance spend for the central path and, after the prior all-S\&M assumption was explained, for the levers, the companion playbook and the engine; all and none of it buying growth as bounds\\
Data exclusion (2 November, 14:50; 22:44) & Remove one registrant from all data, code and replication files; it is outside the estimation sample, so no reported figure changes (Section~\ref{app:data} of the Supplementary Appendix). The stored digest that identifies it can be inverted; the author chose to keep it (22:44)\\
Code and data (2 November, 19:17) & The code and data of all three papers in an open GitHub repository under the author's account, private until this paper is submitted\\
Outside reviews (2 November, 23:23) & Adopt every fix the verification of the three outside reviews supported, among them: ``under (H)'' wherever the decision weight is stated; the acquirer case corrected to its filings, which join the replication package; Nei, Ruus and Indov (2026) cited; the family-dependence of the Holm results and the Monte Carlo error of the Romano--Wolf $p$-value reported, with \BootB{} draws kept; the stress test added, with its conditional claims so stated; third-party texts kept out of the review archives\\
Round-3 tests (3 November, 06:24) & An audit of firm-years against the filings, a specification curve of the empirical claims, a reproduction from a fresh clone, business descriptions of every core-sample registrant for rule R2, and a primary-source check of the legal and tax citations; the retrievals from sec.gov ran in the author's browser\\
Compute share (3 November, 06:24 and 06:34) & The floor of the compute-share range stated in revenue units (the cost per completed task times the labor share) and one range for the compute share at 2025 prices in the text, the calibration and the stress test; the author asked what the correct fix was, the lead agent proposed the two points and she confirmed them, and the lead agent reran the stress test over that range\\
Rule R2 and capitalized software (3 November, 07:06) & Capitalized software read from the filings where company-specific XBRL elements hid it (author's choice); asked how R2 should be used, the author received the lead agent's answer, applied here: the rule as the protocol defined it in the main sample, the two independent readings as a sensitivity reported beside every size they change\\
Stress-test precision (3 November, 12:26 to 14:48) & Asked why the stress test used 10{,}000 draws, the author had the runs redesigned for more precision (a million draws per seed, a larger corner set, search budget, Sobol sample and accountability sub-box), run on her own computer, with the Monte Carlo error reported beside each share, percentile and index\\
Formal verification (3 November, 09:27) & Asked whether the proofs were checked in a proof assistant (they were not), the author chose to have the static model formalized in Lean~4 with Mathlib (Section~\ref{supp:formal} of the Supplementary Appendix); she approved obtaining the Lean toolchain through her computer, which the session's workspace could not reach\\
Further analyses (3 November, about 18:15; 4 November, about 18:00) & Approved the retrieval of public float from the SEC's XBRL frames for a valuation test (3 November); after three analyses had been built and checked for a plain-language companion page, chose to add all three to the Supplementary Appendix before the outside reviewers see the paper again: how valuations weight growth and margin, the cost of not automating and what sets the payback (Section~\ref{supp:extensions} of the Supplementary Appendix)\\
Outside re-reviews (5 November, about 09:30) & Adopted every fix that the verification of the GPT-6.1 Sol and GLM 5.3 re-reviews of the 4 November build classed as valid or partly valid, except two optional ones (an addition to an earlier row of this table, and a sentence on the beta-factor model that awaits a reading of its source), in the wording the issue logs propose, with three choices of her own: \autoref{rem:owncompute} rebuilt on a stated cash basis from Epoch AI's cost tables, with the firm that owns only the servers as the central case and the whole data center also reported; the valuation intervals of Section~\ref{supp:ext-valuation} taken over the bootstrap samples with a positive margin coefficient, with the share of the others in a column of their own; and the floor of the compute-share range relabeled a low scenario value, with no number changed. The Mistral Large 4 re-review, which the lead agent verified, asked for no change to the paper\\
After the re-reviews (5 November, 10:21 to 13:07) & Both GPT reviews ran on GPT-6.1 Sol, not GPT-6.2 as recorded; the round's branch to be merged into the main one; an AI-contributions note on the title page, naming the Claude models and the four outside reviewers, with no AI model listed as an author (13:07); the introduction's new clause on (H), from the DeepSeek re-review, kept; the abstract rewritten in short, plain sentences for business readers, its conditions kept and some details left to the text\\
Rights (2 November, 19:17 for internal use) & Copyright to the author; CC BY-NC-ND 4.0 for the paper; PolyForm Noncommercial 1.0.0 or, at the user's option, PolyForm Internal Use 1.0.0 for the software, so companies may run it internally; no other commercial use without written permission. The licenses cover the text and the software, not the ideas: firms may apply the ideas, while other commercial use of the text or the engine needs the author's written permission\\
\end{longtable}
}
\noindent Table~\ref{tab:decisions} lists the author's decisions. Within them the author delegated methodological choices to the lead agent: prepackaged software (SIC 7372) as the sample; items as first reported; free cash flow after SBC as the economically relevant margin; the statistics and estimands of Section~\ref{sec:empirics}; the range $\alpha\in[0.25,0.45]$; and the scenario values of Tables~\ref{tab:ceiling}, \ref{tab:fullcost} and \ref{tab:calib}.

\subsection{Choices made after results were seen}\label{app:protocol-choices}
Some choices that shape the results were made after some results were known. They are researcher degrees of freedom, not errors, and we list each with its time and reason.
\begin{enumerate}[leftmargin=*,itemsep=0pt,topsep=2pt]
\item \emph{Primary sample for the model-facing results (11:19).} After the author asked what academic rigor requires on reporting results with and without acquisition-affected years, the lead agent made the organic frontier, rather than all firm-years, the primary sample for the model-facing results of the first version, so that the estimand matched the theory, whose growth is bought with expensed spend, and the calibration, which already excluded acquired growth. The exclusion-only robustness check was known at the time: it gave upper-decile slopes of \KnownSlopeOpExcl{} (GAAP) and \KnownSlopeCashExcl{} (after SBC), against \KnownSlopeOpEone{} and \KnownSlopeCashEone{} on all firm-years, and changes in the 95th percentile of the score of \KnownTopOpExcl{} and \KnownTopCashExcl{} points, against \KnownTopOpEone{} and \KnownTopCashEone.
\item \emph{Two pro forma rules (11:44).} A test run printed its organic-sample estimates at 11:41. At 11:44 two rules were added: the pro forma adjustment applies only to acquisition-affected firm-years, and a firm-year whose additive and cumulative readings of the disclosure differ by more than one point is treated as not identified. Without the two rules, the organic upper-decile slope was \TestvTwoLocOp{} on the GAAP basis and \TestvTwoLocCash{} after SBC on the first version's sample, against \TestvTwoRefLocOp{} and \TestvTwoRefLocCash{} with them; on the audited sample the primary slope is \TestPsiOp{} and \TestPsiCash{} without them, against \TestRefPsiOp{} and \TestRefPsiCash{} (Table~\ref{tab:robust}).
\item \emph{Estimand policy (D-3, 13:53).} E1 became the baseline, and a model-facing claim stands only if it holds under E1, E2 and both imputation bounds. The decision followed the round-1 econometrics referees' finding that the differences between the estimands reported in the first version were not robust and reflected which firm-years the organic sample dropped.
\item \emph{Frontier object (D-2, 13:53).} The author asked the lead agent for the right call. The lead agent set the theory's object, $\Psi_\chi$, as primary for the model-facing claims and the calibration of the growth technology, with the local conditional quantile reported as a labeled descriptive statistic. The choice was made before any estimate of $\Psi_\chi$ on the audited sample had been computed, but not blind to its value: the slopes of $\Psi_{0.90}$ on the first version's sample, $-0.20$ on the GAAP basis and $-0.16$ after SBC under E1, were known from the round-1 econometrics report, and they are within 0.02 of the audited estimates. The choice followed that referee's recommendation and weakened the paper's slope claim. Test runs of the superseded local-quantile pipeline on the audited sample, with a tenth of the replications, had printed point estimates and a familywise result in the sample-audit agent's log by 13:39. The sample audit's results notes open with a dated record of the decisions, written at 14:05, before the full estimation run.
\item \emph{Growth weight (D-1, 13:53).} The statement that firms act on a weight on growth of 7 to 13 was withdrawn after eight reports disputed it; the bounds each method gives are reported instead (Section~\ref{sec:empirics}).
\item \emph{Parameters (D-4, 13:53).} Every value of the second assumptions round was adopted in one commit, decided before any calibration with the new values had been run.
\item \emph{Sample rules.} The five rules of Section~\ref{sec:data-rules} respond to findings of the review panels and were fixed before any estimate on the audited sample. Rules R3 (on XBRL data), R4(c) and R5 were applied in a rerun after the missing SEC data had been retrieved; the rerun left every existing estimate unchanged except two flag-threshold variants, which each lost one firm-year. The cheap-acquisition criterion of rule R4 was chosen among five variants after the firm-years each would flag had been listed. Rule R2, set in the protocol with a pre-set keyword test, could not be applied to the main sample until the business descriptions of every core-sample registrant were retrieved on 3 November (until then the \VRtwoFirms{} firms named in the round-1 reports were classified as a sensitivity only); it was then applied to the main sample, together with the correction of capitalized software read from the filings, after the estimates on the earlier sample were known. The re-estimation removed \VRtwoExclFirms{} firms (\VRtwoExclN{} firm-years) and changed the headline values: the firms from 292 to \VFirmsEone, the slope $\Psi_{0.90}$ from $-0.22$ to \PsiOpEone{} (GAAP) and from $-0.16$ to \PsiCashEone{} (cash), the fall of the cash score's 95th percentile from 12.7 to \TopFallCashEone{} points (the GAAP fall of \TopFallOpEone{} is unchanged), the calibration anchor from 13 to \AnchorN{} firm-years, the central 2036 score from 59 to \CenRSix{} and the Romano--Wolf $p$-value of the GAAP decline from 0.097 to \HolmTopRW.
\item \emph{Remote threshold of the accountability floor (second version).} The first draft of the second version used $\bar\pi=1\%$ a year as the central threshold of \autoref{prop:amin}. After the example showed that no all-human design passes at 1\%, the lead agent moved the central value to 5\%, with 1\% kept as a sensitivity, in all three papers; the reason given was that a threshold at which the reference all-human design fails would contradict most filers obtaining effective opinions. The choice was made after the designs at both thresholds were known, and Table~\ref{tab:decisions} records no author decision on it. In the third version neither value is central; both are reported.
\item \emph{Basis of selling and marketing (2 November, 10:15 and 12:35).} The author chose the netted basis for the central path with its 2036 score on each basis before her. A propagation agent then found that automating the growth engine outranks the G\&A cut (Section~\ref{sec:levers-position}) at \lpNGEFirstNet{} of \lpNPositions{} positions on that basis, \lpNGEFirstAllSM{} if all selling and marketing buys growth and \lpNGEFirstNoSM{} if none does; after the prior all-S\&M assumption was explained to her, she chose the netted basis for the levers too, reversing the second version's lever finding.
\item \emph{Romano--Wolf draws (2 November, 23:23).} With the Romano--Wolf $p$-value of the GAAP decline then close to 0.10, the author kept \BootB{} draws rather than rerunning with more; the text reports its Monte Carlo standard error. On the sample re-estimated on 3 November (rule R2 and the capitalized-software correction) the $p$-value is \HolmTopRW, with a standard error of \HolmTopRWSe{} (Section~\ref{sec:empirics}).
\end{enumerate}
The first version also stated that three of its conclusions changed when the primary sample was switched to the organic frontier: a steeper GAAP slope, and a cash-basis size gradient and steepening that did not survive the removal of acquired growth. These statements are withdrawn (Section~\ref{sec:empirics-robust}). The lead agent had reported them to the author as findings before the referees showed that they were not robust.

\subsection{Safeguards and errors caught}
Six safeguards address the known failure modes of AI-generated research: verification passes for every citation; symbolic and numerical checks of the derivations, referees who re-derived them and, for the static model, proofs machine-checked in Lean~4; spot checks of the panel against filed figures and a replication package; numbers in the text generated from the analysis outputs; independent critics and an assumptions register against confirmation bias; and referees on a second model. Section~\ref{supp:protocol} of the Supplementary Appendix describes them and lists the errors that each caught (Table~\ref{tab:errors}).

\subsection{What the protocol does not establish}
Independent agents built on the same underlying model may share blind spots, so their agreement is weaker evidence than agreement among human experts would be. Running referees on a second model reduces this risk but does not remove it, because all the models come from one developer. The machine check covers the static model of Sections~\ref{sec:model-setup} to~\ref{sec:model-statics} only, and it checks that the proofs prove the formal statements; that those statements say what the paper says rests on two audits of the formalization by an agent that had not written it (Section~\ref{supp:formal} of the Supplementary Appendix). The dynamic, valuation, industry, entry and automation results rest on written proofs, symbolic and numerical checks and the referees. Several parameters remain scenario values. The CLAIMS checks and the generated numbers test the text against the pipeline's own outputs, not the inputs against the world. Regulatory and accounting statements were read at primary sources by agents, some only in secondary descriptions (Table~\ref{tab:param-sources}), and no human lawyer, auditor or specialist in accounting standards reviewed them. The statements that we found no prior study of a kind (Sections~\ref{sec:lit} and~\ref{sec:intro}) rest on web searches by the literature agents (27 to 29 October 2026) and the outside reviewers (29 October to 3 November), with each candidate checked at Crossref, RePEc or the publisher; no log of databases and queries was kept, so they are bounded by those searches. The final checks were the author's reading of the manuscript and four independent reviews she commissioned from other large language models, run outside this workflow: DeepSeek V4 Pro (high effort), GPT-6.1 Sol (extra-high effort), GLM 5.3 (high effort) and Mistral Large 4 (high effort). Their findings were verified item by item and answered (Section~\ref{supp:protocol-outside} of the Supplementary Appendix and the response to the outside reviews). Because these models come from other developers, the reviews lower the risk of blind spots shared with the drafting agents. None found a false theorem or a wrong entry in a table; together they found an unjustified approximation in the reading of the cross-sectional slope, an incomplete expansion in a proof, a misdescribed shape of the worked decision weight, mislabeled units and tax bases, source descriptions that overstated what the sources measure, statements about the acquirer's filings that the filings contradict, and an unstated trim in a within-firm specification, which this version corrects. Two kinds of check followed: audits of each revision by an agent on \texttt{claude-fable-5-1} that had not taken part in it, and the five tests described in Section~\ref{app:protocol-rounds}. These found what the reviews had not: the floor of the compute-share range had been set in the wrong units, capitalized software was missing for \CapswKeptFirms{} firms that report it only under company-specific elements, and records of the revision described as done a test that had not finished; each is corrected. The specification curve showed that the empirical magnitudes depend on analysis choices, which the text now reports as ranges. The outside reviews are still reviews by language models, not human peer review: each ran on a single model family, one of them (GPT-6.1 Sol) could not reach the repository and ran none of the code, and none is a human specialist's check of the tax, regulatory or accounting statements. Responsibility for any remaining error is the author's.
